% Polished Dataset Construction section (proofread by Claude). Original text is in main-original.tex.
\section{Dataset Construction}
\label{sec:dataset}
\subsection{Parametric Voxel Generation}
To crystallize broad geometric priors into explicit supervision, GPT-6.1 Sol assists in authoring analytic parametric generator programs for 21 object categories, each partitioned into 40 structural families (47--48 instances per family in each 40,000-model shard). For every instance, a planning module samples independent design parameters, which govern category-specific components, proportions, openings, and support structures. It then deterministically rasterizes the geometry onto a $48^3$ grid ($Y{=}0$ is ground contact) and derives the L1--L5 captions from the same underlying plan. Instead of applying random cell deletion, cavity filling, or post-hoc bridges, density adjustments regenerate the entire shape from its continuous parameters until automated checks confirm 750--1,250 occupied voxels forming a single six-neighbor connected component. Exact and canonical occupancy hashes over the 48 axis-permutation and sign transforms confirm zero cross-shard duplicates, and the geometry plan and version identifiers are stored with each record (Figure~\ref{fig:data-workflow}, where \texttt{shape(params)} denotes the analytic shape and \texttt{emit(x,y,z)} records occupied voxel coordinates).

BrickForma pairs each shape with five bilingual descriptions. L1 names the object and its primary parts, and L2--L5 progressively add structural, proportional, appearance, and viewpoint details about the same source geometry (Appendix~\ref{app:dataset} gives a birdhouse example).

\subsection{Voxel-to-Brick Conversion and Checks}
Converting a connected voxel volume into a buildable brick target requires separating mere spatial adjacency from vertical stud connectivity. Our pipeline tiles each layer with same-color rectangular bricks and analyzes the resulting stud-connected components. When ungrounded fragments arise, local refinement modifies the layout and logs every voxel edit against the source shape, raising the geometry-pass count from 65,503 to 79,298 (13,795 recovered failures). Across the corpus, refinement adds 228,582 voxels and removes 125, and 31,532 passing models require no edits at all.

For challenging layouts, conversion falls back to layer-wise exact-cover search or to a mixed-integer linear program (MILP) that explicitly enforces a stud path to the ground. A minimum-remaining-values (MRV) heuristic branches on the uncovered voxel with the fewest valid candidate bricks. When greedy placement failed the geometric audit, the initial release also used bounded beam search, ranking partial layouts by uncovered voxels and remaining stud components.

Release-level verification confirms that the target voxels form a single six-neighbor component, that every brick component is grounded, and that the brick assembly exactly covers the derived target. Release-level checks allow several independently grounded components; inference-time refinement enforces a single component. Separately, a self-weight screen evaluates the final assembly and, where feasible, each bottom-up prefix of completed layers. It records the tipping margin and estimates the per-stud tensile demand when compression alone cannot balance gravity (Appendix~\ref{app:dataset}).

\subsection{Integrated Release and Training Corpus}
The full refined release passes 79,298 models and leaves 702 geometry failures. A revised release of the batch-B piano models passes all 1,905 piano slots, compared with 1,902 in the full refined release, so the integrated release contains $79{,}298-1{,}902+1{,}905=79{,}301$ passing models. Table~\ref{tab:training-data} summarizes the release counts and the fine-tuning corpus.
\begin{table}[t]
\centering
\small
\setlength{\tabcolsep}{4pt}
\begin{tabular}{lr}
\toprule
\multicolumn{2}{l}{\emph{Dataset (release reports)}} \\
\midrule
Source models / categories / grid & 80,000 / 21 / $48^3$ \\
%Bilingual 
Caption records & 400,000 \\
Failing in full refined release & 702 (0.88\%) \\
Passing in full refined release & 79,298 \\
Piano models restored by revised release & 3 \\
Integrated passing models & 79,301 \\
\midrule
\multicolumn{2}{l}{\emph{Fine-tuning corpus}} \\
\midrule
%396,610 ?
Voxel-model \& caption samples & 396,610 \\ % Training manifest confirmation pending
Training / %validation 
test samples & 356,955 / 39,655 \\
%xiao
% Tokens over 3 epochs & 2,228,044,818 \\ % Training manifest confirmation pending
\bottomrule
\end{tabular}
\caption{Dataset and fine-tuning corpus. The two blocks are counted separately: the upper block comes from the dataset release reports, and the lower block is as recorded for the fine-tuning run.}
\label{tab:training-data}
\end{table}

